\chapter{Signal Serving and Parameter Management}\label{ch:pl:signal-serving-and-parameter-management}

A risk parameter entered as 50 meant fifty basis points to the person who typed it and fifty percent to the strategy that read it. Nothing
checked the unit; the change went to every strategy at once; the first sign was the fill rate. A parameter is code that nobody compiles:
it changes behaviour as much as a release does, it is edited by more people, and it usually has no types, no tests and no history. This
chapter gives parameters and research signals the machinery code already has --- types, review, history, gradual rollout --- and measures,
on ten simulated strategies, what each piece of that machinery saves when fifty means the wrong thing.

\section{From research signal to production input}

A research signal becomes a production input when a strategy reads it on a schedule and trades on it. Between the two sit the questions a
research notebook never asks: which version is live, how old the value may be before it is not to be used, what the strategy does when it is
missing, and who changed it. Book~12 (chapters 24 and 25) answered them for machine-learning features and models --- the feature store,
feature freshness, the model registry and its audit trail. This chapter applies the same answers to the plainer inputs every strategy has:
signals computed by research jobs and parameters typed by people.

\section{The signal service}

\begin{definition}[Signal service]\label{def:pl:signal-serving-and-parameter-management:service}
A \emph{signal service}\index{signal service} serves research signals to production strategies: each value with its version and publication
time, served only while it is fresher than a stated bound, and replaced by a declared fallback when it is stale or missing, with the status
reported to the reader.
\end{definition}

The fallback is part of the signal's contract, chosen when the signal is designed, not by the strategy's author at three in the morning: for
a directional signal, usually zero (no view); for a risk input, usually the most conservative value.

\omcode{code/firm/paramstore/firm_paramstore.py}{148}{165}{The signal service: the latest value if it is fresh, else the declared
fallback, always with the version and a status.}\label{lst:pl:signal-serving-and-parameter-management:service}

With a freshness bound of five minutes and a fallback of zero, a momentum signal published at time 0 and again at 60~seconds is served as
0.7 (version~3, fresh) at 200~seconds and as 0 with the status ``stale-fallback'' at 600~seconds; a signal never published is served as the
fallback with the status ``missing-fallback''. The strategy knows which it got.

\section{Configuration as data}

\begin{definition}[Parameter store, configuration as data]\label{def:pl:signal-serving-and-parameter-management:store}
A \emph{parameter store}\index{parameter store} holds every production parameter of every strategy, each value with the time from which it
applies and the time it was recorded, the change that set it and its approvals. \emph{Configuration as data}\index{configuration as data} is
the practice of keeping configuration in such a store, as typed, versioned records that tools can validate, compare and query, rather than as
text edited in place on production machines.
\end{definition}

A value's two times are the valid time and knowledge time of Book~7 (chapter~3): the time from which it applies and the time the store learned
it. They let the store answer the question every post-mortem asks --- what was live at 10:31 last Tuesday --- in both of its senses: as the
firm knew it then, and as it is known now that a correction has been recorded.

\section{Schemas, approvals and history}

\begin{definition}[Configuration schema]\label{def:pl:signal-serving-and-parameter-management:schema}
A \emph{configuration schema}\index{configuration schema} declares, for each parameter, its canonical unit and the units it may be entered
in, its bounds, the largest change allowed in one step, and the size of change that needs a second approver; every proposed value is converted
and checked against it before anyone approves it.
\end{definition}

The chapter's parameter is an inventory limit, as a share of the strategy's capital: canonical unit a fraction, entered in basis points,
percent or as a fraction, bounded by 2\%, changed by at most 50 basis points at a time, and needing two approvers for any change above 20 basis
points --- the four-eyes principle of One Quant Book~16 (chapter~12) applied to a number.

\omcode{code/firm/paramstore/firm_paramstore.py}{75}{96}{A proposal: converted to the canonical unit, checked against bounds and the
step limit, and told how many approvals it needs; every proposal, refused or not, goes to the audit trail.}\label{lst:pl:signal-serving-and-parameter-management:propose}

The fifty of the hook, entered without a unit, is refused: a value with no unit is not accepted. Entered as ``50 percent'', it is refused twice over: 0.5 is
outside $[0, 0.02]$, and a step of 0.495 exceeds the allowed 0.005. Entered as ``50 bp'' it is accepted as 0.005. A raise to 75 basis points
(a step of 25) waits after one approval, cannot be approved by its author, and applies after the second. Every proposal, approval and
application is a line of Book~12's hash-chained audit trail (\texttt{firm.exptrack}): editing any line breaks the chain, and
\texttt{verify} says where.

\omcode{code/firm/paramstore/firm_paramstore.py}{98}{120}{Approval by someone other than the author, as many as the change needs, and the
bitemporal query: the value in force at a valid time, as known at a knowledge time.}\label{lst:pl:signal-serving-and-parameter-management:approve}

\begin{omfigure}
\begin{tikzpicture}[font=\scriptsize,box/.style={draw,rounded corners=2pt,align=center,minimum height=0.9cm,minimum width=1.8cm,
  fill=omDef!10},>=Stealth]
\node[box] (p) at (0,0) {proposal\\ value + unit};
\node[box] (s) at (2.3,0) {schema\\ units, bounds};
\node[box] (a) at (4.6,0) {approvals\\ by others};
\node[box] (f) at (6.9,0) {flag\\ one strategy};
\node[box] (r) at (9.2,0) {rollout\\ all strategies};
\foreach \x/\y in {p/s,s/a,a/f,f/r} {\draw[->] (\x) -- (\y);}
\node[box,fill=omThm!10,minimum width=9.6cm] (h) at (4.6,-1.5) {bitemporal history (valid and knowledge time) and hash-chained audit trail};
\foreach \x in {p,s,a,f,r} {\draw[->,gray] (\x.south) -- (\x.south|-h.north);}
\node[omThm,font=\scriptsize,anchor=south] at (2.3,0.55) {refused: 50 with no unit, 50 percent};
\node[font=\scriptsize,anchor=south] at (8.05,0.55) {monitor fires: flag off, value rolled back};
\end{tikzpicture}

\omcaption{The life of a parameter change in the store: typed on entry, checked by the schema, approved by others, rolled out behind a flag
to one strategy and then to all, every step recorded with its times in a history that cannot be edited
unnoticed.}\label{fig:pl:signal-serving-and-parameter-management:life}
\end{omfigure}

\begin{example}[What was live at 10:31 last Tuesday]\label{ex:pl:signal-serving-and-parameter-management:asof}
The raise to 75 basis points was recorded on Tuesday at 09:43 and took effect at 10:00. On Thursday a correction was recorded: from Tuesday
10:30 the limit should have been 60. Asked on Wednesday what was live at 10:31 on Tuesday, the store answers 75 basis points; asked now, it
answers 60. Both answers are right, and a post-mortem needs both: what the strategy used, and what it should have used.
\end{example}

\section{Rolling out a change}

\begin{definition}[Feature flag]\label{def:pl:signal-serving-and-parameter-management:flag}
A \emph{feature flag}\index{feature flag} is a named switch in configuration that turns a behaviour or a new value on for chosen strategies,
users or a share of traffic without a new release, so that a change can be rolled out to one first and withdrawn by switching the flag off.
\end{definition}

\begin{definition}[Configuration drift]\label{def:pl:signal-serving-and-parameter-management:drift}
\emph{Configuration drift}\index{configuration drift} is a difference between the configuration a system is declared to run and the one it
actually runs: a server not updated, a value edited in place, a flag left from an old release.
\end{definition}

Flags are how the canary deployment and staged rollout of Book~7 (chapter~21) apply to parameters: the new value is behind a flag switched on
for one strategy, watched, and extended. Drift is what a flag, left in place, becomes: the chapter's \texttt{drift} compares, strategy by
strategy, the declared value with the one each strategy reports it is running, and on the demonstration finds the one strategy running 0.5
where the store says 0.005. The best-known case of both is Knight Capital's: on 1 August 2012 a flag once used for an old function was
repurposed for new code, the new code had not reached one of eight servers, and orders carrying the flag triggered the old code there; the
loss exceeded \$460 million (Book~7, chapter~21, and Book~13, chapter~22, tell it in full).

\section{Fifty of what?}

The chapter's incident runs ten quoting strategies (chapter~11's quoter, one lot a side) for an hour each on Book~10's simulator, each on its
own \texttt{firm.tape} session, with \$4 million of capital each. The intended limit, 50 basis points of capital, is two lots at \$100; read as
50 percent it is two hundred. The change arrives ten minutes into the session. A monitor compares each strategy's fills in every five-minute
window with its usual count for that window --- here, the same session without the change, a better baseline than any real monitor has --- and
fires above 1.3 times; a firing rolls the parameter back. Exposure is the inventory beyond what the intended limit allows (three lots, with
one order pending), in dollars.

\omcode{code/platforms/16-signal-serving-and-parameter-management/python/pl_paramstore.py}{62}{68}{The monitor: the end of the first
five-minute window after the change in which fills exceed 1.3 times the usual count.}\label{lst:pl:signal-serving-and-parameter-management:monitor}

\begin{omfigure}
\begin{tikzpicture}
\begin{axis}[width=12cm,height=5.8cm,xlabel={minutes after the change},ylabel={excess inventory (\$ thousands)},xmin=-1,xmax=20,ymin=0,
  grid=major,grid style={gray!20},tick label style={font=\small},label style={font=\small},table/col sep=comma]
\addplot[color=omThm,very thick,const plot] table[x=minutes,y=all_k]{figdata/platforms/16-signal-serving-and-parameter-management/timeline.csv};
\addplot[color=omDef,thick,const plot] table[x=minutes,y=staged_k]{figdata/platforms/16-signal-serving-and-parameter-management/timeline.csv};
\node[omThm,font=\scriptsize,anchor=west] at (axis cs:5.4,380) {all at once: detected at 5 minutes};
\node[omDef,font=\scriptsize,anchor=south west] at (axis cs:6,45) {staged: the canary alone, detected at 15 minutes};
\node[font=\scriptsize,anchor=east] at (axis cs:19.6,250) {schema check: refused at entry, nothing};
\end{axis}
\end{tikzpicture}

\omcaption{Inventory beyond the intended limit, summed over the strategies that received the wrong value, from the change until the monitor
fires and the value is rolled back, on a ten-second grid. All at once, ten strategies accrue it for five minutes; staged, one strategy for
fifteen; with a schema, the value never enters. Data: \texttt{fig\_paramstore.py}.}\label{fig:pl:signal-serving-and-parameter-management:timeline}
\end{omfigure}

\begin{table}[ht]
\centering\small
\begin{tabular}{lrrr}
\toprule
policy & detection (minutes) & peak excess (\$k) & excess \$-minutes (k) \\
\midrule
all at once & 5 & 410 & 1\,348 \\
staged (one canary for 30 min) & 15 & 30 & 53 \\
schema check & 0 (refused) & 0 & 0 \\
\bottomrule
\end{tabular}
\caption{The units error under three policies: minutes from the change to detection, the peak of the excess inventory summed over affected
strategies, and its integral over time.}\label{tab:pl:signal-serving-and-parameter-management:policies}
\end{table}

\Cref{tab:pl:signal-serving-and-parameter-management:policies} is the result. All at once is detected soonest --- ten monitors watch instead of
one, and the first fires at five minutes --- and costs most: \$410\,000 of excess inventory at the peak and 1.35 million dollar-minutes. The
canary's own monitor takes fifteen minutes to fire (its session is quiet), but only the canary holds the wrong value: \$30\,000 at the peak and
53\,000 dollar-minutes, 25 times less. The schema check costs nothing, because the error never becomes a value. The strategies differ widely in
how soon their fill rate shows the error: five of the ten after five minutes, three after fifteen, one after thirty and one after fifty.

\begin{omfigure}
\begin{tikzpicture}
\begin{axis}[width=11cm,height=5cm,ybar,bar width=12pt,ylabel={minutes to detection},ymin=0,ymax=58,
  symbolic x coords={q01,q02,q03,q04,q05,q06,q07,q08,q09,q10},xtick=data,enlarge x limits=0.06,
  tick label style={font=\small},label style={font=\small},grid=major,grid style={gray!20},table/col sep=comma,
  nodes near coords,nodes near coords style={font=\scriptsize}]
\addplot[fill=omDef!45,draw=omDef] table[x=strategy,y=detect_min]{figdata/platforms/16-signal-serving-and-parameter-management/detection.csv};
\end{axis}
\end{tikzpicture}

\omcaption{Minutes from the units error to the fill-rate monitor's first firing, strategy by strategy, each on its own one-hour session. How soon
an error shows depends on the market the strategy happens to trade: five strategies show it in the first window, one only at the end of the
hour. The staged rollout's canary is q01. Data: \texttt{fig\_paramstore.py}.}\label{fig:pl:signal-serving-and-parameter-management:detection}
\end{omfigure}

\begin{remark}[Detection is not protection]\label{rem:pl:signal-serving-and-parameter-management:detect}
Faster detection did not make all at once safer: it had ten chances to notice and ten strategies accruing exposure while it did. A staged
rollout limits the damage to the canary whatever the monitor's speed; a schema prevents the damage whatever the rollout. The monitor here is
optimistic --- it knows each strategy's counterfactual fill count --- and a real one would be slower, which widens every gap.
\end{remark}

\begin{dated}{2026-09}{Configuration at scale}\label{dat:pl:signal-serving-and-parameter-management:scale}
Facebook's configuration-management paper (Tang and others, SOSP 2015) described a system handling thousands of online configuration changes
a day and trillions of configuration checks, used for product rollouts, experiments, load balancing and model deployment. The SEC's order on
Knight Capital (Release 34-70694, 2013) records the repurposed flag, the server not updated, about 45 minutes of unintended orders and a loss
of more than \$460 million.
\end{dated}

\section{Tutorial: a units error, three ways}\label{tut:pl:signal-serving-and-parameter-management:tut}

\begin{tutorial}
\textbf{Goal.} Serve a signal with a freshness bound, put a parameter under a schema with approvals and history, and measure what a units error
costs under three policies. \textbf{End state:} \cref{fig:pl:signal-serving-and-parameter-management:timeline} and
\cref{tab:pl:signal-serving-and-parameter-management:policies}.
\begin{tutsteps}
\item \textbf{Signals}: \texttt{pl\_paramstore.signals()}.
\item \textbf{Parameters}: \texttt{governance(dir)}: the refused entries, the two approvals, the bitemporal answers, flags, drift, the audit.
\item \textbf{Incident}: \texttt{incident()} runs twenty hours of simulated trading (ten strategies, with and without the error).
\item \textbf{Policies}: read the detection times and exposures from the result.
\end{tutsteps}
\textbf{What to change next.} Lower the monitor's threshold to 1.2 and count its false alarms on the unchanged runs; stage the rollout in three
steps (one, three, all) instead of two.
\end{tutorial}

\section{Build: the parameter store}\label{bld:pl:signal-serving-and-parameter-management:paramstore}

\begin{build}
\bfield{Purpose} Every production input a strategy reads --- signal or parameter --- typed, versioned, approved, auditable and rolled out
gradually.
\bfield{Interface} \texttt{Schema}; \texttt{ParamStore(root)} with \texttt{register}, \texttt{propose}, \texttt{approve}, \texttt{as\_of},
\texttt{history}, \texttt{audit}; \texttt{Flags}; \texttt{drift}; \texttt{SignalService(max\_age, fallback)} with \texttt{publish} and
\texttt{get}.
\bfield{Rules} A value enters with its unit or not at all; the author never approves; changes above a threshold need two approvers; values
carry valid and knowledge times; the audit trail is hash-chained; signals carry a version, a freshness bound and a declared fallback.
\bfield{Acceptance tests} \texttt{code/firm/paramstore/tests/}: units, bounds and steps; approvals (author refused, two distinct approvers);
bitemporal queries and history; the audit trail's verification and its detection of an edited line; flags, drift and the \omterm{def:pl:signal-serving-and-parameter-management:service}{signal service}'s
three statuses.
\bfield{Stretch} Percentage rollouts by hash of the strategy name; automatic rollback wired to the monitor; drift checks against each strategy's
own report at start-up.
\end{build}

\begin{omsources}
\item C.~Tang et al., ``Holistic configuration management at Facebook'', SOSP 2015.
\item US Securities and Exchange Commission, \emph{In the Matter of Knight Capital Americas LLC}, Release No.\ 34-70694, 2013.
\item One Quant Book~7, chapters 3 and 21; One Quant Book~12, chapters 24 and 25; One Quant Book~16, chapter~12.
\end{omsources}

\section{Exercises}

\begin{exercise}[$\star$]\label{exo:pl:signal-serving-and-parameter-management:1}
How many lots is an inventory limit of 50 basis points of \$4 million at \$100 a share, and of 50 percent?
\end{exercise}

\begin{exercise}[$\star$]\label{exo:pl:signal-serving-and-parameter-management:2}
Why does the raise from 50 to 75 basis points need two approvers under the chapter's schema, and the correction to 60 only one?
\end{exercise}

\begin{exercise}[$\star$]\label{exo:pl:signal-serving-and-parameter-management:3}
What does the \omterm{def:pl:signal-serving-and-parameter-management:service}{signal service} return at 400 seconds in the chapter's example, and why?
\end{exercise}

\begin{exercise}[$\star\star$]\label{exo:pl:signal-serving-and-parameter-management:4}
Why was the all-at-once rollout detected sooner than the staged one, and why is it still worse?
\end{exercise}

\begin{exercise}[$\star\star$]\label{exo:pl:signal-serving-and-parameter-management:5}
A post-mortem asks what limit a strategy used at 10:31 on Tuesday. Which of the store's two answers does it want, and when does it want the
other?
\end{exercise}

\begin{exercise}[$\star\star$]\label{exo:pl:signal-serving-and-parameter-management:6}
Which of the chapter's controls would have stopped the Knight Capital incident as the SEC describes it, and which would not?
\end{exercise}

\begin{exercise}[$\star\star\star$]\label{exo:pl:signal-serving-and-parameter-management:7}
\emph{Coding.} Rerun \texttt{incident} with the canary chosen as the strategy whose monitor fires last (q10). What happens to the staged
rollout's detection time and exposure?
\end{exercise}

\begin{exercise}[$\star\star\star$]\label{exo:pl:signal-serving-and-parameter-management:8}
\emph{Find the flaw.} ``Our parameters live in a YAML file in the repository, so they are versioned, reviewed and safe.''
\end{exercise}

\section{Problem: Fifty of What?}

\begin{problem}[{Weekend problem --- a units error under three policies}]\label{pb:pl:signal-serving-and-parameter-management:1}
The ten strategies, the schema and the monitor of the chapter.

\medskip\noindent\textbf{Part I --- The store.}
\begin{enumerate}
\item What does the schema of the inventory limit declare?
\item What happens to 50 entered without a unit, as percent, and as basis points?
\item Who may approve a change, and how many approvals does it need?
\item What are a value's two times, and what question do they answer?
\item What does the audit trail guarantee?
\end{enumerate}

\medskip\noindent\textbf{Part II --- Signals and flags.}
\begin{enumerate}[resume]
\item What does the \omterm{def:pl:signal-serving-and-parameter-management:service}{signal service} return, and with what status, when a signal is stale?
\item Who chooses a signal's fallback, and when?
\item What is a \omterm{def:pl:signal-serving-and-parameter-management:flag}{feature flag} for, and what does it become if left in place?
\item What does the drift check find in the demonstration?
\item What did a repurposed flag do at Knight Capital?
\end{enumerate}

\medskip\noindent\textbf{Part III --- The incident.}
\begin{enumerate}[resume]
\item What do 50 basis points and 50 percent mean in lots?
\item How does the monitor decide, and why is it optimistic?
\item When does each strategy's monitor fire?
\item What are the exposure and time to detection all at once?
\item And staged?
\end{enumerate}

\medskip\noindent\textbf{Part IV --- The verdict.}
\begin{enumerate}[resume]
\item State the \emph{named result}: the exposure accrued before detection by the units error under an all-at-once change, a staged rollout
and a schema check, and the minutes to detection in each case.
\item Why is detection not protection?
\item Which control would you build first?
\item What would you add to catch a units error that stays within bounds?
\item In one sentence: what does a parameter need that a variable does not?
\end{enumerate}
\end{problem}

\section{Interview questions}

\begin{interviewq}[$\star$ \iqroles{developer}]\label{iq:pl:signal-serving-and-parameter-management:1}
How do you stop a units error in a production parameter?
\end{interviewq}

\begin{interviewq}[$\star\star$ \iqroles{developer, researcher}]\label{iq:pl:signal-serving-and-parameter-management:2}
A research signal feeds live strategies. What happens when the job that computes it fails?
\end{interviewq}

\begin{interviewq}[$\star\star$ \iqroles{developer}]\label{iq:pl:signal-serving-and-parameter-management:3}
How would you roll out a change to a risk limit across a hundred strategies?
\end{interviewq}

\begin{interviewq}[$\star\star$ \iqroles{developer}]\label{iq:pl:signal-serving-and-parameter-management:4}
What is \omterm{def:pl:signal-serving-and-parameter-management:drift}{configuration drift}, and how would you detect it?
\end{interviewq}

\begin{interviewq}[$\star\star\star$ \iqroles{developer}]\label{iq:pl:signal-serving-and-parameter-management:5}
Design a parameter management system for a trading firm.
\end{interviewq}
